% Generated draft: requires mathematical and source review.
\documentclass[11pt]{article}
\usepackage[margin=25mm]{geometry}
\usepackage{fontspec}
\setmainfont{Latin Modern Roman}
\usepackage{amsmath,amssymb,mathtools,unicode-math}
\setmathfont{Latin Modern Math}
\usepackage{enumitem,needspace,xurl,hyperref}
\setlength{\parindent}{0pt}
\setlength{\parskip}{5pt}
\setlength{\emergencystretch}{3em}
\newcommand{\N}{\mathbb{N}}
\newcommand{\Z}{\mathbb{Z}}
\newcommand{\Q}{\mathbb{Q}}
\newcommand{\R}{\mathbb{R}}
\newcommand{\C}{\mathbb{C}}
\newcommand{\F}{\mathbb{F}}
\newcommand{\A}{\mathbb{A}}
\newcommand{\PP}{\mathbb{P}}
\newcommand{\E}{\mathbb{E}}
\newcommand{\eps}{\varepsilon}
\newcommand{\dd}{\,\mathrm{d}}
\begin{document}
\section*{UnsolvedMath ID: 30004847}
\textbf{Source problem code:} catalogue entry
\subsection{Definitions and mathematical statement}
For terminology, relative to a chosen commutative ground ring \(k\), a differential graded cocommutative coalgebra consists of a chain complex \(C\) of \(k\)-modules with a coassociative chain map \(\Delta:C\to C\otimes_k C\) satisfying \(\tau\circ\Delta=\Delta\), where \(\tau(x\otimes y)=(-1)^{|x||y|}y\otimes x\) for homogeneous elements. A differential graded Lie algebra over \(k\) is a chain complex with a degree-zero graded Lie bracket for which the differential is a graded derivation. The source does not specify the ground ring, grading convention, or additional coalgebra conditions such as counitality or conilpotence; these definitions do not select those missing hypotheses.

\textbf{Source question.} Can \(I\)-chains be used to construct models of spaces in the settings of differential graded cocommutative coalgebras and differential graded Lie algebras?

\textbf{Unresolved formulation.} The supplied record does not define \(I\)-chains or their indexing category, specify the class of spaces, or define the required meaning of a model. In particular, it does not state the comparison functors, equivalences, or recovery properties that such models must satisfy, nor whether both algebraic settings must be related by a single construction. Thus the question is preserved as a research target; the supplied information does not determine a fully quantified existence statement.

The supplied dated assessment says that homotopy colimits of \(I\)-chains produce \(E_\infty\)-coalgebras and that a known obstruction to the expected left-induced model structures prevents a complete Lie or cocommutative-coalgebra model of the requested kind. The record does not identify the precise obstruction or its hypotheses. This assessment is supplied curation, not an independently verified theorem or a claim that every possible construction is impossible.
\subsection{Short English statement}
The question asks whether \(I\)-chains can provide algebraic models of spaces using differential graded cocommutative coalgebras and Lie algebras. The intended spaces, coefficients, construction, and notion of equivalence remain unspecified. The supplied open-status label and obstruction assessment require independent review.
\subsection{Sources}
\begin{itemize}
\item Local UnsolvedMath registry: permanent ID \(30004847\), source problem code \texttt{catalogue entry}, in the Oberwolfach Reports Open Problems collection. The supplied record labels the problem open and its statement \texttt{corrected\_verified}; these are registry annotations, not independent mathematical verification.
\item Supplied source citation: \emph{Homotopical Algebra and Higher Structures} (2022), Oberwolfach Reports, pp. 2497--2560, DOI \texttt{10.4171/owr/2021/46}. Literature-source label: Oberwolfach Report 46. \url{https://doi.org/10.4171/owr/2021/46}. The record gives proposed year 2021.
\item The original statement appends the text ``Canonical \(E_\infty\)-operads involved in homotopy colimits of \(I\)-chain complexes'' and the name Birgit Richter. This preserves the supplied attribution context; the separate proposer field is empty.
\item Supplied literature source: \emph{Comonoidal properties of I-chains}. \url{https://www.math.uni-hamburg.de/en/personen/richter/ichains100822.pdf}.
\item Literature assessment checked at \texttt{2026-08-22}: the supplied curation labels the question open and reports the \(E_\infty\)-coalgebra construction and obstruction described above. This is supplied dated literature triage derived from the OWR curation, not independently verified source consultation or a current literature assessment.
\end{itemize}
\end{document}
